% =============================================================================
%                              MODULE OVERVIEW
% =============================================================================
\section*{Module Overview}
\addcontentsline{toc}{section}{Module Overview}

\subsection*{About This Module}

This module provides an overview of the various challenges that underpin \textbf{robot autonomy} and \textbf{AI-enabled robotics}. It covers the distinct technical problems involved in allowing robots to perform tasks autonomously and provides students with the algorithms used to solve them.

\textbf{Module code:} 6CCSAPRA \hfill \textbf{Semester:} 2 (AY 2025/26)

\textbf{Module leaders:} Dr.\ Khen Elimelech \& Dr.\ Matteo Leonetti

\bigseparator

\subsection*{Module Aims}

The module aims to:
\begin{itemize}
  \item Introduce \textbf{robot kinematics}---the fundamental language for modelling robot movement.
  \item Cover \textbf{motion planning}---the algorithmic framework enabling complex robots to move to desired positions, even in challenging environments.
  \item Examine \textbf{object grasping} and \textbf{task and motion planning (TAMP)} for interactive, complex tasks.
  \item Explore \textbf{localization and mapping}, \textbf{handling uncertainty}, and how to leverage \textbf{machine learning} in robotics.
\end{itemize}

\bigseparator

\subsection*{Learning Outcomes}

By the end of this module, you will be able to:

\begin{center}
\renewcommand{\arraystretch}{1.4}
\begin{tabular}{@{} c p{12cm} @{}}
\toprule
\textbf{Code} & \textbf{Learning Outcome} \\
\midrule
PRA1 & \textbf{Identify} the main problems of interest at the intersection of AI, ML, and Robotics that underpin robot autonomy, such as planning, perception, state estimation, and learning. \\
PRA2 & \textbf{Apply, compare, and analyse} solution methods and algorithms for those problems. \\
PRA3 & \textbf{Design and integrate} the algorithmic components needed to create an autonomous system capable of performing a given task. \\
\bottomrule
\end{tabular}
\end{center}

\bigseparator

\subsection*{Key Topics Overview}

The module is structured around the following core areas:

\begin{center}
\renewcommand{\arraystretch}{1.3}
\begin{tabular}{@{} l p{10cm} @{}}
\toprule
\textbf{Topic} & \textbf{Description} \\
\midrule
Robot Kinematics & Mathematical models describing robot motion; forward/inverse kinematics; degrees of freedom. \\
Motion Planning & Path and trajectory planning; sampling-based algorithms (RRT, PRM); configuration space. \\
Object Grasping & Grasp planning; force closure; manipulation strategies. \\
Task \& Motion Planning & Combining symbolic task planning with geometric motion planning. \\
Localization \& Mapping & SLAM; probabilistic state estimation; sensor fusion. \\
Handling Uncertainty & Probabilistic robotics; belief states; Bayesian filtering. \\
Machine Learning & Learning from demonstration; reinforcement learning for robotics. \\
\bottomrule
\end{tabular}
\end{center}

\bigseparator

\subsection*{What Examiners Like to Ask}

Based on the learning outcomes, expect questions such as:
\begin{itemize}
  \item \textbf{Conceptual:} ``Explain the difference between forward and inverse kinematics.''
  \item \textbf{Algorithmic:} ``Describe how RRT works and analyse its completeness properties.''
  \item \textbf{Comparative:} ``Compare PRM and RRT for motion planning.''
  \item \textbf{Applied:} ``Given a robot arm configuration, compute the end-effector position.''
  \item \textbf{Design:} ``Design an autonomous system for a pick-and-place task.''
\end{itemize}

\TODO\ --- Add specific exam question patterns once past papers are available.
